\section{Manipulación de puntuaciones: técnicas prácticas de manipulación de puntajes}

El primer enfoque consiste en manipular directamente la función de puntuación durante el proceso de eliminación de ruido. Este tipo de métodos no cuenta con una explicación teórica estricta, pero suelen funcionar bien en la práctica, habiendo generado numerosos casos de éxito en los campos de imagen y video.

\subsection{RePaint: un caso clásico de reparación de imágenes}

\textbf{RePaint} es un ejemplo típico de manipulación de puntuaciones; resuelve el problema de \textbf{reparación de imágenes (Inpainting)}: dado una imagen parcialmente oculta, rellenar las zonas faltantes para que el resultado sea natural y coherente con el fondo conocido.

\begin{importantbox}{La idea central de RePaint}
En cada paso de eliminación de ruido, se divide el punto de datos intermedio $x_t$ en dos partes para procesarlas:
\begin{itemize}
    \item \textbf{Zona de primer plano} (zona por rellenar): ejecutar el \textbf{proceso inverso} (reverse process) para generar contenido gradualmente a partir del ruido.
    \item \textbf{Zona de fondo} (zona conocida): ejecutar el \textbf{proceso hacia adelante} (forward process), añadiendo ruido a la imagen limpia conocida hasta el nivel de ruido correspondiente al paso temporal $t-1$.
\end{itemize}
Luego se unen ambas partes para obtener $x_{t-1}$, repitiéndose este proceso hasta que $t=0$.
\end{importantbox}

El flujo del algoritmo específico es el siguiente:
\begin{enumerate}
    \item Dado el punto de datos intermedio actual $x_t$, ejecutar un paso de proceso inverso para obtener la zona de primer plano $x_{t-1}^{\text{fg}}$.
    \item Ejecutar un paso de proceso hacia adelante en la imagen de fondo limpia conocida hasta el paso temporal $t-1$, obteniendo $x_{t-1}^{\text{bg}}$.
    \item Combinar ambos utilizando una máscara (mask): tomar $x_{t-1}^{\text{fg}}$ para el primer plano y $x_{t-1}^{\text{bg}}$ para el fondo.
    \item Repetir el proceso anterior.
\end{enumerate}

\begin{warningbox}{RePaint carece de garantías teóricas}
Por qué se conserva el fondo es intuitivo: en cada paso, añadimos ruido de nuevo a partir de la imagen limpia. Sin embargo, no existe garantía teórica sobre la \textbf{integración perfecta} (seamless integration) entre primer plano y fondo. En uso real, efectivamente existen casos fallidos donde el primer plano y el fondo no coinciden. El docente enfatiza: ``No hay teoría aquí, solo ideas prácticas.''
\end{warningbox}

\subsection{Manipulación de mapas de atención}

Además de manipular directamente la puntuación/ruido, también se puede lograr una guía mediante la manipulación de los \textbf{mapas de atención} (Attention Maps) dentro de la red predictora de ruido, incluyendo:

\begin{itemize}
    \item \textbf{Manipulación de auto-atención}: modificar los pesos de atención en el módulo de atención automática para influir en las asociaciones entre diferentes regiones de la imagen.
    \item \textbf{Manipulación de atención cruzada}: modificar la correspondencia entre los tokens de texto y las regiones de imagen en la atención cruzada, controlando así la posición y las atributos de objetos específicos.
\end{itemize}

\begin{knowledgebox}{Ejemplos de aplicación de la manipulación de atención cruzada}
Mediante la manipulación de la atención cruzada, se pueden lograr las siguientes tareas (todas sin necesidad de volver a entrenar el modelo de difusión):
\begin{itemize}
    \item \textbf{Transferencia de estilo}: extraer información estructural de una imagen y información de estilo/apariencia de otra para generar una nueva imagen.
    \item \textbf{Control de disposición}: dado un texto descriptivo y cajas delimitadoras 2D, controlar la posición de cada objeto en la imagen. Por ejemplo, especificar que el ``coche amarillo'' esté dentro del recuadro amarillo a la derecha y el ``castillo'' dentro del recuadro verde a la izquierda.
    \item \textbf{Edición de video}: mediante inversión (inversion) y manipulación de atención, mover objetos en un video a nuevas posiciones mientras se mantiene la naturalidad del video.
\end{itemize}
\end{knowledgebox}

\subsection{Ventajas y desventajas de la manipulación de puntuaciones}

\begin{table}[H]
\centering
\begin{tabular}{p{0.45\textwidth}p{0.45\textwidth}}
\toprule
\textbf{Ventajas} & \textbf{Desventajas} \\
\midrule
Directo, intuitivo y fácil de implementar & Sin garantías teóricas \\
Numerosos casos de éxito en los campos de imagen/video & Existen casos fallidos; los resultados son impredecibles \\
No requiere entrenamiento adicional ni datos & Mayormente métodos ad-hoc diseñados para tareas específicas \\
Extensible a otras modalidades de datos & Difícil de generalizar a nuevos tipos de guía \\
\bottomrule
\end{tabular}
\caption{Comparación de ventajas y desventajas de los métodos de manipulación de puntuaciones}
\end{table}

\subsection{Resumen de la sección}

La manipulación de puntuaciones es una clase de técnicas de guía en tiempo de inferencia altamente prácticas. Su idea central consiste en modificar directamente la función de puntuación o los mapas de atención en cada paso del proceso de eliminación de ruido, haciendo que la trayectoria de generación se incline hacia la dirección deseada. Aunque carece de garantías teóricas, ya existen numerosos casos de éxito en tareas como reparación de imágenes, transferencia de estilo, control de disposición y edición de video. A continuación discutiremos una clase de métodos más sistemáticos y con base teórica: la guía basada en problemas inversos.

